当区间的加点操作易于实现，而删点操作难以实现时，使用此种算法。

\paragraph{询问排序}
首先对序列进行分块，块大小通常取 $\sqrt{n}$。然后将所有询问按下述规则排序：
\begin{itemize}
    \item 以左端点 $l$ 所在的块编号为第一关键字，升序排序。
    \item 以右端点 $r$ 为第二关键字，升序排序。
\end{itemize}

\paragraph{指针移动}
我们按块的顺序处理所有询问。
\begin{enumerate}
    \item \textbf{同块内询问}：如果一个询问的左右端点 $l, r$ 在同一个块内，由于块长不超过 $\sqrt{n}$，直接暴力扫描该区间计算答案即可。
    \item \textbf{跨块询问}：对于左端点在块 $T$ 内的所有询问，我们进行如下操作：
    \begin{itemize}
        \item 首先将莫队区间的左端点 $l$ 初始化为块 $T$ 的右边界加一 ($R[T]+1$)，右端点 $r$ 初始化为 $R[T]$ 。此时区间为空。
        \item \textbf{移动右指针}：由于同一块内的询问是按 $r$ 升序排列的，我们可以只向右移动右指针 $r$，不断加点，直到满足当前询问的 $r$ 。对于整个块 $T$ 的所有询问，指针 $r$ 只会单调向右移动。
        \item \textbf{移动左指针}：对于每一个询问，我们从 $R[T]+1$ 出发，将左指针 $l$ 向左移动到询问所需的位置，这个过程只涉及加点操作。移动距离不会超过块的大小。
        \item \textbf{回答与回滚}：回答当前询问后，我们需要撤销刚才左指针的所有移动，让 $l$ 回滚到 $R[T]+1$ 的位置，以恢复基准状态，准备下一个询问。这个回滚操作只撤销数据结构的变化，而不需要更新答案。
    \end{itemize}
    \item 处理完一个块内的所有询问后，按照相同方式处理下一个块。
\end{enumerate}

\paragraph{只删不加}
当删点操作易于实现，而加点操作困难时使用。其流程与只加不减对称。
\begin{itemize}
    \item \textbf{排序}：左端点所在块升序，右端点\textbf{降序}。
    \item \textbf{指针移动}：对于左端点在块 $T$ 内的询问，将莫队区间初始化为 $[L[T], n]$ 。右指针 $r$ 只会单调向左移动（删点），左指针 $l$ 从 $L[T]$ 临时向右移动（删点），回答后再回滚。
\end{itemize}
